% !TeX root = ../教材题库_全部条目_转写工作稿.tex
% Source: XJ-B1-C06-RV-P264-Y1051
% 第6章 / 章末复习 复习题六 / 教材 p.264 / PDF 第271页
% OCR未校对：公式与文字均可能有误；未完成数学TeX转写；禁止直接发给学生
% 原题及相邻原解截图（核对用）：../题目原图/XJ-B1-C06-RV-P264-Y1051-01.webp
% 本文件只在转写完成后作为题干使用；OCR 条目尚不是可靠题干。
\begin{bankitem}{复习题六 第4题}{第6章 / 章末复习 复习题六 / 教材 p.264 / PDF 第271页}{XJ-B1-C06-RV-P264-Y1051}{OCR未校对：公式与文字均可能有误；未完成数学TeX转写；禁止直接发给学生}
\begin{ocrdraft}
[教材p.264 / PDF第271页]
\par
4．在一次测量中得到的A样本数据为：82，84，84，86，86，86，88，88，
\par
88，8\&若B样本数据恰好是，A样本数据每个都加2后所得数据，则A，B两样
\par
本的下列数字特征对应相同的是（）
\par
（B)平均数
\par
（A)众数
\par
(C)中位数
\par
(D)标准差
\end{ocrdraft}
\end{bankitem}
